%% =================================================================
%% Chen et al. Measurement 264 (2026) 120300.
%% Reconstruction of page 6 (printed p. 6).
%% Contents: Fig. 8 (applied-force study at 0.5/1.5/2.5/4 MPa),
%% prose covering Section 3.1 end, Eq. (1) tactile-softness
%% definition, Section 3.2 Tactile softness of materials discussion.
%% =================================================================

\documentclass[11pt]{article}

\usepackage[margin=0.85in]{geometry}
\usepackage{amsmath, amssymb}
\usepackage{mathrsfs}
\usepackage{xcolor}
\usepackage{fancyhdr}

\pagestyle{fancy}
\fancyhf{}
\lhead{\small Z.~Chen et al.}
\rhead{\small Measurement 264 (2026) 120300}
\cfoot{\small 6 $\to$ \arabic{page}}
\renewcommand{\headrulewidth}{0.4pt}

\setcounter{page}{1}
\setcounter{equation}{0}
\setlength{\parskip}{6pt}
\setlength{\parindent}{0pt}

\begin{document}

% ---------- Fig. 8 placeholder ----------
\begin{center}
\fbox{\begin{minipage}[c][6.5cm][c]{0.95\textwidth}
\centering
\textbf{Fig. 8.}\ Effect of the applied force on the tactile
softness of materials.\\[4pt]
\textcolor{gray}{\footnotesize
3D imaging of compression-depth distribution when the tactile
tomography system scanned the hierarchical structure model (Fig.~7)
with a soft surface layer under a pressure of
\textbf{(a)} 0.5 MPa, \textbf{(b)} 1.5 MPa, \textbf{(c)} 2.5 MPa,
and \textbf{(d)} 4 MPa respectively. Each panel is a rainbow-colored
rendered 3D surface (Z scale 0--6 mm). At low force (0.5 MPa) only
the uppermost staircase plateau resolves; as force increases more
layers of the stepped interior become visible, until at 4 MPa the
full staircase structure is recovered.]}
\end{minipage}}
\end{center}

\vspace{0.4em}

\textbf{Fig. 8.}\ Effect of the applied force on the tactile softness
of materials. 3D imaging of compression depth distribution when the
tactile tomography system scanned the hierarchical structure model
with a soft surface layer under a pressure of (a) 0.5, (b) 1.5,
(c) 2.5, and (d) 4 MPa, respectively.

\vspace{0.6em}

% ---------- Prose (Section 3.1 tail) ----------
\pdfcomment[icon=Comment,color=cyan]{User input: How does the Chen paper recognize the softness?\textbackslash \textbackslash Claude comments: The excerpt gives the complete mechanism of softness recognition in one self-contained passage, moving from physical setup to formal definition in six steps: \textbackslash \textbackslash \textbackslash \textbackslash  (1) PHYSICAL MECHANISM: "applying a constant force through a tip     pressing against the soft material" — the measurement is     force-in, deformation-out, with the constant-force regime     enforced by the closed-loop feedback control system. \textbackslash \textbackslash \textbackslash \textbackslash  (2) CONTROL ARCHITECTURE: "This capability is attributed to the     closed-loop feedback control system" — the control system     (described in Section 2.3 on page 5) is what makes the     constant-force regime possible. Pressure sensor reads tip     force; PLC compares to preset threshold; rodless cylinder     stops Z-axis descent when the threshold is met. \textbackslash \textbackslash \textbackslash \textbackslash  (3) CALIBRATION SUBSTRATES: "a series of two-component epoxy     adhesives (AB glues) with Shore hardnesses of 3, 7.5, 12.5,     17.5, 25, 30, 35, 50, and 60 HC" — 9 calibrated samples     covering a 20× hardness range, measured independently by a     commercial Shore tester (LX-C) to ground-truth their     hardnesses. \textbackslash \textbackslash \textbackslash \textbackslash  (4) OBSERVED RELATIONSHIP: "the AB glues with different Shore     hardnesses exhibit different compression depths. The average     compression depth decreases as the Shore hardness of the AB     glues increases... demonstrating a negative correlation" —     harder sample \$\textbackslash Rightarrow\$ smaller compression depth. This     monotonic relation is the empirical basis for using     compression depth as a proxy for softness. The quantitative     fit (Fig. 5d) is \$Y = 6.23 - 1.62\textbackslash ,X\^\{\}1 + 1.15\textbackslash ,X\^\{\}2\$ with     \$R\^\{\}2 = 0.99596\$. \textbackslash \textbackslash \textbackslash \textbackslash  (5) DEFINITIONAL MOVE: "the softness of soft materials can be     characterized by the compression depth" — the recognition     primitive is \$\textbackslash Delta d\$ (compression depth) itself. \textbackslash \textbackslash \textbackslash \textbackslash  (6) FORMAL DEFINITION (Eq. 1): Softness \$= \textbackslash Delta d / P = \textbackslash Delta     d \textbackslash cdot S / F\$, where \$\textbackslash Delta d\$ is compression depth, \$F\$     is applied force, \$S\$ is contact area, and \$P = F/S\$ is     pressure. This normalizes compression depth by pressure so     softness values are comparable across different probe forces     and contact geometries. \textbackslash \textbackslash \textbackslash \textbackslash  To answer the question in one sentence: Chen recognizes softness by applying a known constant force (set by the closed-loop feedback controller) through a 0.5-mm-diameter probe tip onto the sample, measuring the resulting compression depth \$\textbackslash Delta d\$ at that point (via the grating sensor on the Z-axis), and computing Softness \$= \textbackslash Delta d \textbackslash cdot S / F\$. This is a ... [truncated in sticky; see journal for full]}\sethlcolor{hlcyan}\hl{recognize softness by applying a constant force through a tip
pressing against the soft material. This capability is attributed
to the closed-loop feedback control system. As shown in Fig.~5a, a
series of two-component epoxy adhesives (AB glues) with Shore
hardnesses of 3, 7.5, 12.5, 17.5, 25, 30, 35, 50, and 60 HC were
prepared with a size of $15 \times 15$ mm$^2$ and a thickness of
6 mm, and their hardness values were measured by a Shore hardness
tester (LX-C) (Fig.~5b). The tactile tomography system pressed each
AB glue sample with a constant pressure of 2.5 MPa. The scanning
point for each AB glue is 12, within a scanning area of
$12 \times 12$ mm$^2$. The scanning results are shown in Fig.~5c,
where the AB glues with different Shore hardnesses exhibit different
compression depths. The average compression depth decreases as the
Shore hardness of the AB glues increases (as shown in Fig.~5d),
demonstrating a negative correlation. This indicates that the
softness of soft materials can be characterized by the compression
depth. Therefore, based on the tactile tomography system, the
relationship between the softness and compression depth can be
defined as}
\begin{equation}
\text{Softness} \;=\; \frac{\Delta d}{P}
  \;=\; \frac{\Delta d \cdot S}{F}
\end{equation}
where $\Delta d$ is the compression depth, $F$ is the force when
the tip pressing the sample, and $S$ is the contact area between
the tip and the sample.

\subsection*{3.2.\ Tactile softness of materials}

The size and thickness of the sample must be taken into account when
testing Shore hardness. The contact position of the specimen should
be at least 12 mm from the edge (for Shore A, B, C, and D), and the
thickness must not be less than 6 mm. Consequently, the compression
depth will vary if the size and thickness of the sample are less
than 12 mm and 6 mm, respectively. Similar to tactile softness [38],
which is perceived through the mechanoreceptors in human skin, the
tactile softness of materials is proposed for use in a tactile
tomography system to simulate human haptic perception to recognize
the softness of soft materials, taking into consideration factors
such as sample size, sample thickness, and the applied force. As
shown in Fig.~6a, a frame model was fabricated by 3D printing with
acrylonitrile butadiene styrene (ABS). The AB glue samples were
filled in the frame model to form a series of samples with different
sizes and the same thickness of 6 mm, in which the length and width
were decreased from 8.0 mm to 3.2 mm (8.0, 6.4, 5.1, 4.0, and 3.2
mm, respectively) (Fig.~6b). Subsequently, the samples were scanned
by the tactile tomography system under a constant pressure of 4 MPa.
Fig.~6c presents that the sample with a length \& width of 8 mm \&
8 mm possesses the largest compression depth, and the compression
depth decreases with decreasing length or width of the sample. The
compression depth of the central area of each sample are larger than
that of the edge area (Fig.~6d). These results demonstrate that the
tactile tomography system can effectively simulate human haptic
perception, allowing for the differentiation of tactile softness
distribution in materials when the size of the soft matter is less
than 12 mm.

In addition, a hierarchical structure model was fabricated by 3D
printing with ABS, and the height difference between each adjacent
stage is about 0.5 mm (as shown in Fig.~7a). Then the AB glue was
filled in the hierarchical structure model to form a soft surface
layer with different thickness (Fig.~7b). The image reconstructed by
the projection data from the tactile tomography system scanning this
sample with a constant pressure of 4 MPa is shown in Fig.~7c, the
compression depth of the samples with different thickness surface
layers exhibits clear gradation. The compression depth difference
between each adjacent thickness is also about 0.5 mm (Fig.~7d),
which consistent with the height difference of each adjacent stage.
Significantly, by combining the compression depth and its
corresponding location information ($x$ and $y$), Fig.~7c also
shows the internal hierarchical structure of the sample, indicating
that the tactile tomography system can obtain the internal structure
information of the sample with a soft surface layer based on tactile
softness of materials.

The object that is buried deep cannot be sensed by the fingers
unless a large force is applied. Similarly, the tactile tomography
system is

\end{document}
